\documentclass[11pt]{article}
\usepackage{amsmath, amssymb}
\usepackage[utf8]{inputenc}
\usepackage[margin=1in]{geometry}
\usepackage{graphicx}
\newcommand{\fit}[1]{\resizebox{\ifdim\width>\linewidth\linewidth\else\width\fi}{!}{#1}}
\setlength{\parindent}{0pt}
\begin{document}
\section*{Proof of \texttt{commutative-ring-is-ring}}
\textbf{Claim.}
\begin{align*}
  \forall s.\; (\mathsf{is{-}commutative{-}ring}(s) \Rightarrow \mathsf{is{-}ring}(s))
\end{align*}
\bigskip

\subsection*{Tactics used}
\begin{itemize}
\item \texttt{sp} -- start the proof: install the claim as the initial goal
\item \texttt{di} -- direct inference: break the goal at its top connective -- move an implication's hypothesis into the assumptions, split a conjunction, or introduce a universally quantified variable
\item \texttt{mac-h} -- macete in a hypothesis: unfold a defined predicate inside an assumption
\item \texttt{ai} -- antecedent inference: decompose a structured assumption, named by its number or its formula
\item \texttt{ass} -- assert: close the goal directly from the assumptions
\end{itemize}

\subsection*{Notation}
\begin{itemize}
\item $\Rightarrow$ -- implies
\item $\in$ -- set membership
\item $\wedge$ -- and
\end{itemize}

\subsection*{Derivation}
Each step shows the tactic applied; the goal nodes it opens and which one is now in focus (the bracketed \texttt{[k]} of the workspace display); the focus node's assumptions (\textbf{A1}, \textbf{A2}, \dots{}, the numbering the tactic args refer to); and the goal $G$ it leaves, written $\vdash G$.  Each sequent is an \texttt{align*} environment: break a long formula by hand with \texttt{\textbackslash\textbackslash\ \&} and indent at the alignment column.

\medskip\noindent\textbf{0.}\quad \texttt{sp}\hfill{\small adds node 0 $\;\cdot\;$ focus 0}\\[2pt]
\begin{align*}
  \vdash\quad & \forall s.\; (\mathsf{is{-}commutative{-}ring}(s) \Rightarrow \mathsf{is{-}ring}(s))
\end{align*}
\medskip\noindent\textbf{1.}\quad \texttt{(di)}\hfill{\small adds node 1 $\;\cdot\;$ focus 1}\\[2pt]
\begin{align*}
  \vdash\quad & (\mathsf{is{-}commutative{-}ring}(s) \Rightarrow \mathsf{is{-}ring}(s))
\end{align*}
\medskip\noindent\textbf{2.}\quad \texttt{(di)}\hfill{\small adds node 2 $\;\cdot\;$ focus 2}\\[2pt]
\begin{align*}
  \mathbf{A1.}\quad & \mathsf{is{-}commutative{-}ring}(s) \\
  \vdash\quad & \mathsf{is{-}ring}(s)
\end{align*}
\medskip\noindent\textbf{3.}\quad \texttt{(mac-h is-commutative-ring-def ...)}\hfill{\small adds node 3 $\;\cdot\;$ focus 3}\\[2pt]
\begin{align*}
  \mathbf{A1.}\quad & (\mathsf{is{-}ring}(s) \wedge \forall a \in a(s).\; \forall b \in a(s).\; (\mathit{mul}(s)(a, b) = \mathit{mul}(s)(b, a))) \\
  \vdash\quad & \mathsf{is{-}ring}(s)
\end{align*}
\medskip\noindent\textbf{4.}\quad \texttt{(ai ...)}\hfill{\small adds node 4 $\;\cdot\;$ focus 4}\\[2pt]
\begin{align*}
  \mathbf{A1.}\quad & \forall a \in a(s).\; \forall b \in a(s).\; (\mathit{mul}(s)(a, b) = \mathit{mul}(s)(b, a)) \\
  \mathbf{A2.}\quad & \mathsf{is{-}ring}(s) \\
  \vdash\quad & \mathsf{is{-}ring}(s)
\end{align*}
\medskip\noindent\textbf{5.}\quad \texttt{(ass)}\hfill{\small adds no nodes}\\[2pt]
\emph{QED.}\par

\end{document}
